% from ega0-19-fr.tex
\oldpage[0\textsubscript{IV-1}]{82}

\begin{proof}
\begin{enumerate}
  \item[(i)] Si $C$ est une $A$-algèbre discrète,
  $\varphi:C\to C/\mathfrak J$ l'homomorphisme canonique de $C$ sur une
  $A$-algèbre quotient quelconque de $C$, le \emph{seul} $A$-homomorphisme
  de $A$ dans $C/\mathfrak J$ est l'homomorphisme composé
  $A\xrightarrow{\psi}C\xrightarrow{\varphi}C/\mathfrak J$, où
  $\psi$ est l'homomorphisme définissant la structure de $A$-algèbre de
  $C$~; comme $\psi$ est continu, la condition de (19.3.1) est trivialement
  vérifiée.
  \item[(ii)] Soient $\alpha:A\to B$, $\beta:B\to C$ les homomorphismes
  continus définissant respectivement la structure de $A$-algèbre sur $B$ et
  celle de $B$-algèbre sur $C$, de sorte que $\beta\circ\alpha$ définit la
  structure de $A$-algèbre sur $C$. Soient $E$ une $A$-algèbre discrète,
  $\mathfrak L$ un idéal nilpotent de $E$, $u:C\to E/\mathfrak L$ un
  $A$-homomorphisme continu, de sorte que $u\circ\beta\circ\alpha$ est
  l'homomorphisme définissant la structure de $A$-algèbre de
  $E/\mathfrak L$. Comme $B$ est une $A$-algèbre formellement lisse, le
  $A$-homomorphisme continu $u\circ\beta:B\to E/\mathfrak L$ se
  factorise en $B\xrightarrow{v}E\to E/\mathfrak L$, où $v$ est un
  $A$-homomorphisme continu~; $v$ et $u\circ\beta$ définissent alors sur
  $E$ et $E/\mathfrak L$ respectivement des structures de $B$-algèbre
  topologique, pour lesquelles $E/\mathfrak L$ est encore l'algèbre quotient
  (discrète) de la $B$-algèbre $E$. En outre, $u$ est un $B$-homomorphisme
  continu, donc se factorise en $C\xrightarrow{w}E\to E/\mathfrak L$,
  où $w$ est un $B$-homomorphisme continu~; comme
  $v\circ\alpha$ est le $A$-homomorphisme définissant la structure de
  $A$-algèbre sur $E$, $w$ est bien un $A$-homomorphisme continu, d'où notre
  assertion.
  \item[(iii)] Soient $C$ une $A'$-algèbre topologique discrète,
  $\mathfrak J$ un idéal nilpotent de $C$,
  $u:B\otimes_A A'\to C/\mathfrak J$ un $A'$-homomorphisme continu. Si on
  compose $u$ et l'homomorphisme canonique
  $\rho:B\to B\otimes_A A'$, on obtient (puisque
  $A\to B\xrightarrow{\rho}B\otimes_A A'$ est égal au composé
  $A\to A'\xrightarrow{\sigma}B\otimes_A A'$) un $A$-homomorphisme
  continu qui, par hypothèse, se factorise donc en
  $B\xrightarrow{v}C\to C/\mathfrak J$, où $v$ est un $A$-homomorphisme
  continu. L'égalité des composés $A\to B\xrightarrow{v}C$ et
  $A\to A'\xrightarrow{w}C$ entraîne donc l'existence d'un homomorphisme
  continu d'anneaux $f:B\otimes_A A'\to C$ tel que
  $v=f\circ\rho$ et $w=f\circ\sigma$ ($\mathbf 0_{\rm I}$, 7.7.6)~;
  comme les homomorphismes composés
  $B\xrightarrow{\rho}B\otimes_A A'\xrightarrow{f}C\to C/\mathfrak J$
  et
  $B\xrightarrow{\rho}B\otimes_A A'\xrightarrow{u}C/\mathfrak J$
  (resp.
  $A'\xrightarrow{\sigma}B\otimes_A A'\xrightarrow{f}C\to C/\mathfrak J$
  et
  $A'\xrightarrow{\sigma}B\otimes_A A'\xrightarrow{u}C/\mathfrak J$)
  sont égaux, on a bien la factorisation
  $u:B\otimes_A A'\xrightarrow{f}C\to C/\mathfrak J$, ce qui établit
  (iii).
  \item[(iv)] La topologie considérée sur $S^{-1}A$ (resp. $T^{-1}B$) est
  naturellement celle pour laquelle un système fondamental de voisinages de
  $0$ est formé des $S^{-1}\mathfrak J$ (resp. $T^{-1}\mathfrak R$), où
  $\mathfrak J$ (resp. $\mathfrak R$) parcourt un système fondamental
  d'idéaux ouverts dans $A$ (resp. $B$) ($\mathbf 0_{\rm I}$, 7.6.1). Si
  $\alpha:A\to B$ est l'homomorphisme canonique, il est clair que
  l'homomorphisme canonique $\alpha':S^{-1}A\to T^{-1}B$ déduit de
  $\alpha$ est continu. Soient alors $C$ une $S^{-1}A$-algèbre topologique
  discrète, $\mathfrak J$ un idéal nilpotent de cette algèbre,
  $u:T^{-1}B\to C/\mathfrak J$ un $S^{-1}A$-homomorphisme continu~; alors
  le composé $B\to T^{-1}B\xrightarrow{u}C/\mathfrak J$ est un
  $A$-homomorphisme continu qui, par hypothèse, se factorise en
  $B\xrightarrow{v}C\to C/\mathfrak J$, où $v$ est un $A$-homomorphisme
  continu. Comme pour tout $t\in T$, $t/1$ est inversible dans $T^{-1}B$,
  $u(t/1)$ est inversible dans $C/\mathfrak J$. Comme $\mathfrak J$ est
  nilpotent, tout élément de la classe $u(t/1)$ dans $C$, et en particulier
  $v(t)$, est inversible dans $C$ ($\mathbf 0_{\rm I}$, 7.1.12), et par suite
  $v$ se factorise en $B\to T^{-1}B\xrightarrow{w}C$~; comme $v$ est
  continu, il en est de même de $w$ ($\mathbf 0_{\rm I}$, 7.6.6), et c'est
  un $S^{-1}A$-homomorphisme car le composé
  $A\to B\to T^{-1}B\xrightarrow{w}C$ est égal à
  \[
    A\to S^{-1}A\xrightarrow{\alpha'}T^{-1}B\xrightarrow{w}C,
  \]
  donc $S^{-1}A\xrightarrow{\alpha'}T^{-1}B\xrightarrow{w}C$ est
  l'homomorphisme canonique définissant sur $C$ la structure de
  $S^{-1}A$-algèbre. Enfin, les homomorphismes composés
  $B\to T^{-1}B\xrightarrow{w}C\to C/\mathfrak J$ et
  $B\to T^{-1}B\xrightarrow{u}C/\mathfrak J$ étant égaux, le même
  raisonnement montre que $u$ est bien égal au composé
  $T^{-1}B\xrightarrow{w}C\to C/\mathfrak J$, d'où (iv).
  \item[(v)] Enfin, (v) est immédiat, la donnée d'un $A$-homomorphisme
  continu de $B=\prod_{i=1}^n B_i$ dans $C$ (resp. $C/\mathfrak J$)
  équivalant à celle de $n$ $A$-homomorphismes continus $B_i\to C$
  (resp. $B_i\to C/\mathfrak J$) et tout $A$-homomorphisme continu
  $B_j\to C$ (resp. $B_j\to C/\mathfrak J$) donnant par composition un
  $A$-homomorphisme continu $B\to B_j\to C$
  (resp. $B\to B_j\to C/\mathfrak J$).
\end{enumerate}
\end{proof}

